\section{Explicit Constructions and Distance Scaling}
\label{sec:compact-constructions}\label{sec:allocation-response}
\subsection{TailSpline and the symmetric construction}\label{sec:tailspline-details}\label{sec:bm-construction}
Let $D\epsilon=(\epsilon_{q+1}-\epsilon_q)_{q=1}^{n-1}$ and $H=D^\top D+e_ne_n^\top$. A zero value of $\epsilon^\top H\epsilon$ requires constant increments and a zero terminal increment, hence $\epsilon=0$. Thus $H$ is positive definite. For $n>1$, the proposed $\epsilon_q=c(n+q)(n-q+1)$, with $c=3/[n(n+1)(2n+1)]$, satisfies
\[
\epsilon_1-\epsilon_2=2c,\quad
2\epsilon_q-\epsilon_{q-1}-\epsilon_{q+1}=2c,\quad
2\epsilon_n-\epsilon_{n-1}=2c.
\]
Therefore $H\epsilon=2c\mathbf1$, the equality-constrained stationarity equation. The increments are positive and sum to one; summing them yields Eq.~\eqref{eq:tailspline}. The objective minimum is $2c=6/[n(n+1)(2n+1)]$. For $n=1$, the constraint gives $\epsilon_1=1$, agreeing with the formulas. Positivity makes additional nonnegativity constraints inactive.

Adding $\epsilon_1^2$ yields the symmetric objective, with unique solution
\begin{equation}
\epsilon_q^{\rm BM}=\frac{6q(n-q+1)}{n(n+1)(n+2)},\qquad
m_q^{\rm BM}=\frac{q(q+1)(3n+2-2q)}{n(n+1)(n+2)}.
\label{eq:bm-formula}
\end{equation}
Both profiles are monotone with $m_0=0,m_n=1$.

\paragraph{What the objective smooths.}
On the geometric native grid, Eq.~\eqref{eq:increment-log-gap} makes each change between neighboring log gaps equal to $(\epsilon_{q+1}-\epsilon_q)\log s$. The terminal change is $-\epsilon_n\log s$ because the fully interpolated tail has constant displacement. Thus $(\log s)^2$ times Eq.~\eqref{eq:tailspline-objective} is exactly the squared log-gap variation within the transition and at its low-frequency junction. The high-frequency junction is deliberately unpenalized.

At the terminal gap,
\[
\epsilon_n^{\rm Pro}=\frac2{n+1},\quad
\epsilon_n^{\rm TS}=\frac6{(n+1)(2n+1)},\quad
\frac{\epsilon_n^{\rm TS}}{\epsilon_n^{\rm Pro}}=\frac3{2n+1}.
\]
The opposite entry behavior is also explicit: $\epsilon_1^{\rm Pro}=2/[n(n+1)]$, whereas $\epsilon_1^{\rm TS}=3/(2n+1)$. TailSpline trades a larger entry jump for a smaller terminal jump. A useful boundary prior need not minimize both. In the continuum limit, $m^{\rm TS}(u)=(3u-u^3)/2$ solves $\min\int_0^1(m''(u))^2du$ under $m(0)=0,m(1)=1,m'(1)=0$; the natural entry condition is $m''(0)=0$. The finite-grid formula is used in every model experiment, not this limiting polynomial.

\subsection{An exact equal-displacement control}
\label{sec:tailspline-dose-control}
Use local indices $q=0,\ldots,n$. Write $T_q=m_q^{\rm TS}$, $B_q=m_q^{\rm BM}$, $U_q=q/n$, and $P_q=q(q+1)/[n(n+1)]$. Define the front-loaded profile $F_q=2U_q-P_q$ and the control
\begin{equation}
C_q=(1-w_n)U_q+w_nF_q,\qquad w_n=\frac{3n}{2(2n+1)}.
\label{eq:tailspline-dose-control}
\end{equation}
Substitution of the finite-grid formulas gives $T=(1-w_n)B+w_nF$. Subtracting $C$ gives
\begin{equation}
m^{\rm TS}-m^C=(1-w_n)(m^{\rm BM}-m^{\rm Uni}).
\label{eq:tailspline-shape-contrast}
\end{equation}
In local coordinates, the explicit difference is
\begin{equation}
T_q-C_q=(1-w_n)(B_q-U_q)=\frac{q(n-q)(2q-n)}{2n(n+1)(2n+1)}.
\label{eq:tailspline-antisymmetric-contrast}
\end{equation}
This cumulative-profile difference is antisymmetric under $q\mapsto n-q$, vanishes at both endpoints, and sums to zero. Since $U$ and $F$ have positive increments and $0<w_n<1$, $C$ is monotone with $C_0=0,C_n=1$. Extending both $T$ and $C$ by the same constant outer bands gives identical endpoints and total log displacement at a common scale. Their equal increment centroids follow from Eq.~\eqref{eq:displacement-centroid}, rather than imposing a further constraint. For $n=1,2$, $B=U$ and $T=C$; for wider bands the contrast tests higher-order shape.

For $q<n/2$, $T_q<C_q$; for $q>n/2$, $T_q>C_q$. TailSpline thus
allocates less stretch in the entry half and more in the tail half, at the same
total displacement. The clean T--C results are in Appendix~\ref{sec:clean-shape-new}.

\subsection{Cosh density and finite installation}
\label{sec:proofs}\label{sec:table-identities}
Let $S_\rho(t)=\int_t^1\rho(\phi)d\phi$. For $\alpha>0$, $\beta\ge0$, the
Cosh transport minimizes
\begin{equation}
\Capp[\rho]=\frac\alpha2\int_0^1\rho^2+
\frac\beta2\int_0^1 S_\rho^2,\qquad \rho\ge0,\quad\int_0^1\rho=1.
\label{eq:Capp}
\end{equation}
The first term penalizes concentrated log-frequency density. By Fubini,
$\int_0^1S_\rho(t)^2dt=\mathbb E[\min(\Phi_1,\Phi_2)]$ for independent
$\Phi_i\sim\rho$, so the second penalizes joint occupancy toward the slow
end. These declared structural preferences define the transport objective.
Coercivity and weak lower semicontinuity on the closed convex constraint set
give existence; the first term makes the objective strictly convex. For
$g(\phi)=\int_0^1\min(\phi,\psi)\rho(\psi)d\psi$, stationarity is
$\alpha\rho+\beta g+\lambda=0$. Since $g''=-\rho$, $g'(1)=0$ and
$g'(0)=1$, this gives $\rho''=\tau^2\rho$, $\rho'(0)=-\tau^2$ and
$\rho'(1)=0$, where $\tau^2=\beta/\alpha$. Thus the unique solution is
\begin{equation}
\rho_\tau(\phi)=\frac{\tau\cosh[\tau(1-\phi)]}{\sinh\tau}.
\label{eq:rho-tau-closed}
\end{equation}
It is positive and has unit mass, so it satisfies the inequality constraint.
At $\beta=0$ the solution is uniform. Integrating and inverting its CDF yields
Eq.~\eqref{eq:warp}. Equivalently, for $h=Q'$,
$\Capp=\tfrac12\int_0^1[\alpha/h+\beta(1-u)^2h]du$.

For midpoint quantiles $u_k=(k+1/2)/K$, anchored installation uses
\begin{equation}
z_k=\frac{Q_\tau(u_k)-Q_\tau(u_0)}{Q_\tau(u_{K-1})-Q_\tau(u_0)}
\le\frac{k}{K-1},\qquad \omega_k=\exp(-a-Rz_k).
\label{eq:anchored-cosh-order}
\end{equation}
The inequality follows from convexity of $Q_\tau$ and its endpoint chord.
Unanchored midpoint installation instead uses $\omega_k=b^{-Q_\tau(u_k)}$.
Its uniform limit is midpoint Geo; standard Geo uses $k/K$, differing by the
common frequency factor $b^{-1/(2K)}$. Each experimental protocol below
specifies its actual grid and support convention. Native denotes the inherited
frequency table of a released model.
\subsection{Scale response and finite phase changes}
\label{sec:allocation-scale-response}
Write $\nu_k(s)=\omega_k^N s^{-m_k(s)}$ and let
$A_k(s)=\omega_k^N/\nu_k(s)$ be the wavelength multiplier.
For the index ramp in the released YaRN implementation
\citep{peng2024yarn}, with fixed interior coordinate $0<t<1$,
\begin{equation}
A_Y(s,t)=\frac{1}{1-t+t/s},\qquad
\lim_{s\to\infty}A_Y(s,t)=\frac{1}{1-t}.
\label{eq:yarn-interior-scale-limit}
\end{equation}
A fixed frequency blend $\gamma+(1-\gamma)/s$ similarly has limit
$1/\gamma$ for $\gamma>0$; the fully interpolated tail $\gamma=0$ grows
with $s$. MrPro and TailSpline instead have
$A_P=s^{P_q}$ and $A_T=s^{T_q}$ for their fixed finite-grid exponents.
Their log-scale derivatives are
\begin{equation}
\frac{\partial\log A_Y}{\partial\log s}
=\frac{t}{s(1-t)+t},\qquad
\frac{\partial\log A_P}{\partial\log s}=P_q,\qquad
\frac{\partial\log A_T}{\partial\log s}=T_q.
\label{eq:allocation-scale-derivatives}
\end{equation}
These identities follow by differentiation. They describe constructor
responses to $s$, rather than task-performance guarantees.

Index and turn-count ramps use different blend weights. All numerical
comparisons here use one shared band and the index-ramp equation of the
released YaRN implementation (revision \texttt{995db5b}).
They are frequency-table comparisons, not a re-evaluation of the YaRN
baseline runs in the MrRoPE paper.

\paragraph{Magnitude of the TailSpline intervention.}
Direct subtraction of the finite-grid formulas gives
\begin{equation}
T_q-P_q=\frac{q(n-q)(3n+q+1)}{n(n+1)(2n+1)}\ge0,
\qquad \frac{A_T}{A_P}=s^{T_q-P_q}.
\label{eq:tailspline-wavelength-order}
\end{equation}
For the Llama $n=17$ grid, the maximum wavelength ratio is
$1.328763$, $1.765612$, $2.346080$ and $3.117385$ at $s=2,4,8,16$.
At pair $k=26$ ($q=8$), the multipliers are:
\begin{center}
\begin{tabular}{@{}lrrr@{}}
\toprule
Scale & YaRN index blend & MrPro & TailSpline\\
\midrule
$4$ & $1.545455$ & $1.385674$ & $2.423868$\\
$16$ & $1.789474$ & $1.920093$ & $5.875134$\\
\bottomrule
\end{tabular}
\end{center}
The YaRN/MrPro ordering at an interior pair can therefore change with
scale. TailSpline redistributes substantial wavelength growth throughout
the transition, in addition to its smaller terminal extra-gap jump.

\subsection{Two reference responses and finite phase intervals}
\label{sec:allocation-reference-intervals}
The native reference is $\omega$ and the fully interpolated reference is
$\omega/s$. For a declared phase tolerance $0<\theta\le\pi$, define
\begin{equation}
D_N(\nu;\theta)=\frac{\theta}{|\omega-\nu|},\qquad
D_D(\nu;\theta)=\frac{\theta}{|\nu-\omega/s|},
\label{eq:allocation-reference-distances}
\end{equation}
with infinity at a zero denominator. These are the maximal intervals
starting at zero on which the respective unwrapped phase deviation is at
most $\theta$. Since $\omega/s\le\nu_T\le\nu_P\le\omega$,
\begin{equation}
D_D(\nu_T;\theta)\ge D_D(\nu_P;\theta),\qquad
D_N(\nu_T;\theta)\le D_N(\nu_P;\theta).
\label{eq:allocation-reference-order}
\end{equation}
The exact operator difference is
$2|\sin(d(\nu-\nu_{\rm ref})/2)|$.
Thus the distances quantify two specified phase references; they are not
effective context lengths of a model.

\paragraph{Tail-junction phase residual.}
At $q=n-1$, $m_q=1-\epsilon_n$ and hence
$\nu-\omega/s=(\omega/s)(s^{\epsilon_n}-1)$. For $n\ge2$ and $s>1$,
$0<\epsilon_n^{\rm TS}\le\epsilon_n^{\rm Pro}$. Since
$(e^x-1)/x$ is increasing for $x>0$, dividing the two residuals gives
Eq.~\eqref{eq:tail-phase-ratio}. It quantifies the approach to the stretched
distance reference at the final interior pair.

\paragraph{Dependency stretch and response.}
Let a reference dependency have separation $d_0>0$ and its extended separation
be $\alpha d_0$. The residual phase in channel $q$ is
$\delta_M=\omega d_0(\alpha s^{-m_q^M}-1)$. For an interior pair with
$T_q>P_q$ and $s>1$, the smaller absolute residual switches from P to T at
\begin{equation}
\alpha_q^*=\frac{2}{s^{-T_q}+s^{-P_q}}.
\label{eq:dependency-crossover}
\end{equation}
Indeed, with $v_M=s^{-m_q^M}$,
$\delta_T^2-\delta_P^2=(\omega d_0)^2\alpha(v_T-v_P)
[\alpha(v_T+v_P)-2]$. Thus P has smaller residual for
$0<\alpha<\alpha_q^*$, and T for $\alpha>\alpha_q^*$.
The relevant scale is the dependency's stretch, not the total input length.

To allow a nonzero content phase, write a reference-aligned complex
coefficient as $a+ib$, with $a>0$ and $|b|\le\kappa a$. Its response at residual
$\delta$ is $f(\delta)=a\cos\delta-b\sin\delta$. If two residuals have the
same sign and $0\le r=|\delta_A|<t=|\delta_B|\le\pi$, then
\begin{equation}
f(\delta_A)-f(\delta_B)\ge
2a\sin\frac{t-r}{2}
\left(\sin\frac{t+r}{2}-\kappa\left|\cos\frac{t+r}{2}\right|\right).
\label{eq:finite-coherent-response}
\end{equation}
This follows from the cosine and sine difference identities and
$|b|\le\kappa a$; an allowed endpoint value of $b$ attains the bound.
A positive bracket guarantees a larger response. At $\kappa=0$ it recovers
the aligned cosine comparison; beyond that case the allowable phase offset
is explicit. These are fixed-content responses. Their contribution to
attention depends on competing keys, as specified in
Appendix~\ref{sec:content-competition}. For example, $a=1$, $b=-1$ and
positive residuals $0.1/0.3$ give a response difference of $-0.156$,
illustrating why the phase condition matters.

\paragraph{Equal displacement does not imply a small phase intervention.}
On the same Llama $s=4$ mathematical grid, the largest T/C unwrapped
phase difference over integer distances $0,\ldots,32767$ is
$10.672514$ radians. The corresponding T/P value is $108.197815$.
For finite comparisons one can use
\begin{equation}
\rho_H(\delta)=\frac4H\sum_{d=0}^{H-1}\sin^2(d\delta/2)
=2-2\frac{\sin(H\delta/2)}{H\sin(\delta/2)}
\cos((H-1)\delta/2),
\label{eq:allocation-finite-phase-cost}
\end{equation}
with removable singularities evaluated by continuity. The fixed-slot
rotation RMS is $[K^{-1}\sum_k\rho_H(\nu_k-\widetilde\nu_k)]^{1/2}$.
At $H=32768$, it is $0.448329$ for T/C and $0.687543$ for T/P.
These values compare unscaled rotary kernels, before the shared
cosine/sine gain and the attention normalization.
For $H|\delta|\ll1$,
$\rho_H(\delta)=((H-1)(2H-1)/6)\delta^2+O(H^4\delta^4)$.
The small parameter is therefore distance times frequency difference;
matching total displacement alone does not supply it.

\paragraph{Turn boundaries across native bases.}
\label{sec:turn-normalization}
The $32/1$-turn rule places the transition in a comparable native phase range.
For a geometric grid with $c=\log b/K$, assume both boundaries lie inside the
sampled table without clipping. Let $l$ be the last slot above $32$ turns,
$h$ the first below one, $n=h-l$ and $u=q/n$. Define
$\eta_+=\log(L\omega_l/(64\pi))$ and
$\eta_-=\log(2\pi/(L\omega_h))$. Both belong to $(0,c]$, and
\begin{equation}
nc=\log32+\eta_++\eta_-,\qquad
L\omega_{l+q}=64\pi\,32^{-u}
 \exp\!\big((1-u)\eta_+-u\eta_-\big).
\label{eq:turn-normalization}
\end{equation}
These identities follow from $\omega_{l+q}=\omega_l e^{-cq}$.
Together with $d\nu=(d/L)(L\omega)s^{-m}$, they explain why a larger base
alone does not remove a turn-selected allocation change. The remaining
geometric differences include finite sampling, rounding and transition width;
model task behavior additionally depends on learned content use.
The fine-grid limit is $c\to0$, not increasing the base at fixed $K$.
